\documentclass[../main.tex]{subfiles}
\begin{document}

\section{Rationale}
Small software teams---student groups in particular---usually manage their work with chat
messages, spreadsheets or a general-purpose tool such as Trello or Jira. Each choice has a cost.
Chat loses information; spreadsheets know nothing about status, dependencies or deadlines; and the
large tools are heavy, mouse-driven, require a constant connection and charge for the features a
team actually needs, such as Scrum reports. Meanwhile, two developments make a better tool
practical: browsers can now hold a full working copy of a project locally (the
\emph{local-first} idea~\cite{kleppmann2019localfirst}), and large language models can turn
unstructured text---meeting notes, a vague goal---into structured, actionable data.

This project builds a team project management system around three ideas: every action
is reachable from the keyboard; the board keeps working when the network does not; and AI is used
for the tedious parts of management (summaries, minutes, assignment, task breakdown, sprint
retrospectives) while always returning structured, checkable output.

\section{Objectives}
\begin{itemize}
  \item Build a web application in which users register, create projects and invite members, with
        roles (Project Manager or Member) assigned \emph{per project} rather than per account.
  \item Provide a four-column Kanban board and a table view, fully operable from the keyboard,
        with dependencies between tasks and an automatically computed critical path.
  \item Support the Scrum process: sprint planning, a product backlog, sprint start and completion,
        burndown and velocity charts, and an AI-written retrospective.
  \item Integrate an LLM (Groq, model \code{openai/gpt-oss-120b}) for five AI features, with a
        three-tier fallback so that the application never fails when the AI service is unavailable.
  \item Keep a personal project fully usable offline, while preventing conflicting offline edits in
        shared projects.
  \item Deploy a production-configured instance on a public domain for the team to test, and
        verify the system with automated tests.
\end{itemize}

\section{Scope}
The system targets small teams (roughly 2--15 people) working on software projects. It is a web
application usable on desktop and mobile browsers, in English and Vietnamese. Out of scope are time
tracking, billing, Gantt charts, native mobile applications and integration with third-party issue
trackers.

\section{Methodology}
The team applied the process the product supports. Work was organised in sprints with a shared
backlog; each change was accompanied by tests, and every non-trivial decision was recorded. The
project keeps eight living documents (D1 Requirements, D2 Module map, D3 Architecture, D4 API and
data contracts, D5 Tests and acceptance criteria, D6 Risks and policies, D7 Development book and
decision log, D8 Requirements traceability matrix) in the same repository as the code, so that a
change to the code and the documents describing it are reviewed together. Testing followed a
\emph{characterisation-first} approach: before changing untested code, a test describing its
current behaviour was written; tests themselves were checked by \emph{mutation
testing}~\cite{demillo1978mutation}---deliberately breaking the code and confirming that a test
fails.

\section{Team organisation and task allocation}
The team has three members from class Computer Engineering K23A. Table~\ref{tab:tasks} shows how
the work was divided.

\begin{longtable}{|p{3.2cm}|p{2.2cm}|p{8.2cm}|}
\caption{Task allocation among team members}\label{tab:tasks}\\
\hline
\textbf{Member} & \textbf{Role} & \textbf{Responsibilities} \\
\hline
\endfirsthead
\hline
\textbf{Member} & \textbf{Role} & \textbf{Responsibilities} \\
\hline
\endhead
Phạm Văn Huynh \newline (KTMT K23A) & Team leader \newline Scrum Master &
\begin{itemize}[leftmargin=*,topsep=0pt]
 \item Project planning, sprint organisation and progress tracking
 \item Requirements analysis and the documentation set D1--D8
 \item Scrum module: sprint lifecycle, backlog planning, burndown, velocity, AI retrospective
 \item AI integration with Groq and prompt design; language-aware AI answers
 \item Deployment: Docker, MySQL, Cloudflare Tunnel on \code{koshi.huynhca.com}
 \item Test data, end-to-end smoke tests, demo video, this report
\end{itemize} \\
\hline
Phạm Minh Tú \newline (KTMT K23A) & Lead developer \newline (architecture, frontend) &
\begin{itemize}[leftmargin=*,topsep=0pt]
 \item Overall architecture and technology choices
 \item Vue 3 frontend: Pinia store, Kanban and table views, keyboard engine
 \item Dependency graph, topological sort and critical path
 \item Local-first storage (IndexedDB) and the offline write policy
 \item Landing page, profile page, localisation (English/Vietnamese)
\end{itemize} \\
\hline
Đàm Đức Đôn \newline (KTMT K23A) & Developer \newline (backend, quality) &
\begin{itemize}[leftmargin=*,topsep=0pt]
 \item FastAPI backend: authentication, projects, per-project roles and invitations
 \item Comments, @mentions, replies, attachments, notifications, task history
 \item Database schema, Alembic migrations, migration from SQLite to MySQL
 \item Continuous integration (GitHub Actions), backend and frontend test suites
 \item AI fallback logic and key rotation
\end{itemize} \\
\hline
\end{longtable}

\section{Structure of the report}
Besides the introduction and conclusion, the report has five chapters. Chapter~1 presents the
theoretical background: Kanban, Scrum, dependency graphs and the critical path, local-first
software, LLMs and the technologies used. Chapter~2 analyses the requirements. Chapter~3 describes
the system design: architecture, data model, API, AI cascade, security and offline policy.
Chapter~4 presents the implementation with screenshots, the testing strategy and its results, and
the deployment. Chapter~5 describes the AI agent ``software factory'' used to build the system:
context engineering, the stations of the production line, executable guardrails, unit and
integration testing, and the packaging, operation and monitoring of the product.

\end{document}
